\chapter{Algorithmic and Quantitative Analysis}

This chapter formalizes the major algorithms used in the project and explains
their practical effect. The purpose is not only to present formulas, but also
to connect each formula to implementation decisions in sender and receiver
modules.

\section{Bit Encoding and Message Cost}
In early versions, each character is split into 8 bits and transmitted as
separate UDP datagrams. If a message has length $L$ characters, then the number
of payload datagrams in pure bit mode is
\begin{equation}
N_{\text{bit}} = 8L + 8,
\end{equation}
where the additional $8$ corresponds to transmitting the null terminator as one
more byte.

This immediately explains why bit mode is inefficient. Datagram count grows
linearly with a large constant factor.

\section{Parity-Based Error Detection}
For versions with parity, sender computes parity of one byte:
\begin{equation}
p = b_0 \oplus b_1 \oplus \cdots \oplus b_7,
\end{equation}
where $b_i$ is the $i$th bit of the character and $\oplus$ denotes XOR.

Receiver computes parity again and compares. If mismatch is found, the packet
is marked corrupted. Parity detects all odd-bit errors but cannot detect every
even-bit error pattern. This is why later versions move to checksum.

\section{Stop-and-Wait Service Time}
In stop-and-wait, sender transmits one unit and waits for ACK before next send.
Approximate per-packet cycle time is
\begin{equation}
T_{\text{cycle}} = T_{\text{tx}} + T_{\text{prop}} + T_{\text{proc}} + T_{\text{ack}}.
\end{equation}

Ideal throughput under stop-and-wait is therefore bounded by
\begin{equation}
\text{Throughput}_{\text{SAW}} \approx \frac{\text{payload per packet}}{T_{\text{cycle}}}.
\end{equation}

Even if channel capacity is high, waiting for each ACK keeps effective
throughput low when round-trip delay is non-trivial.

\section{Checksum in TCP-like Versions}
In advanced versions, a one's complement style checksum is computed over key
header fields and payload bytes. Conceptually:
\begin{align}
S &= \sum \text{(header fields + payload bytes)}, \\
S' &= (S \bmod 2^{16}) + \left\lfloor\frac{S}{2^{16}}\right\rfloor, \\
\text{checksum} &= \sim S'.
\end{align}

At the receiver, checksum field is reset, recomputed, and compared. This gives
stronger detection than parity and supports multi-byte payload packets.

\section{Sliding Window Efficiency}
With sender window size $W$, up to $W$ packets can be in flight. In an ideal
pipeline, normalized channel utilization can be approximated as
\begin{equation}
U \approx \min\left(1, \frac{W}{1 + 2a}\right),
\end{equation}
where $a = T_{\text{prop}}/T_{\text{tx}}$.

This relation explains why larger windows increase throughput until the path is
fully utilized.

\section{Chunking Gain}
If payload per packet increases from 1 byte to $C$ bytes (for example,
$C=8$), the number of data packets for message length $L$ becomes
\begin{equation}
N_{\text{chunk}} = \left\lceil\frac{L}{C}\right\rceil.
\end{equation}

Relative packet count reduction compared to one-byte payload is approximately
\begin{equation}
\text{Gain factor} \approx C,
\end{equation}
ignoring control packets and retransmissions.

\clearpage
\section{Adaptive RTO Estimation}
Versions 12 and above use smoothed RTT estimation and variation-based timeout.
Given sample RTT $R$, smoothed RTT ($SRTT$) and RTT variation ($RTTVAR$) are
updated as follows:
\begin{align}
RTTVAR &\leftarrow (1-\beta)\,RTTVAR + \beta\,|SRTT - R|, \\
SRTT   &\leftarrow (1-\alpha)\,SRTT + \alpha\,R.
\end{align}

Then retransmission timeout is estimated by
\begin{equation}
RTO = SRTT + K \cdot RTTVAR,
\end{equation}
and bounded in implementation between predefined minimum and maximum values.

This model is aligned with practical transport engineering guidance and avoids
both aggressive and overly delayed retransmissions.

\section{Congestion Window Dynamics}
In congestion-controlled versions, sender maintains congestion window $cwnd$ and
slow-start threshold $ssthresh$.

\subsection{Slow Start}
When $cwnd < ssthresh$, growth is fast (approximately one MSS-equivalent unit
per ACK, resulting in near exponential growth per RTT in simplified models).

\subsection{Congestion Avoidance}
When $cwnd \ge ssthresh$, growth becomes gradual, approximating additive
increase.

\subsection{Loss Response}
On timeout event:
\begin{align}
ssthresh &\leftarrow \max\left(2, \frac{cwnd}{2}\right), \\
cwnd &\leftarrow 1.
\end{align}

On triple duplicate ACK in fast retransmit logic:
\begin{align}
ssthresh &\leftarrow \max\left(2, \frac{cwnd}{2}\right), \\
cwnd &\leftarrow ssthresh + 3.
\end{align}

After successful recovery, window is deflated near threshold level. This avoids
fully restarting from 1 in many recoverable loss cases.

\section{Latency Components}
Observed transfer completion time includes multiple components:
\begin{equation}
T_{\text{total}} = T_{\text{setup}} + T_{\text{data}} + T_{\text{retrans}} + T_{\text{teardown}}.
\end{equation}

Where:
\begin{enumerate}
    \item $T_{\text{setup}}$ includes handshake steps,
    \item $T_{\text{data}}$ includes normal in-order transfer,
    \item $T_{\text{retrans}}$ includes recovery delay from loss/corruption,
    \item $T_{\text{teardown}}$ includes connection close signaling.
\end{enumerate}

This decomposition helps explain why different versions show different
completion times for the same message length.

\section{Reliability Metrics}
Useful metrics for experiment reporting are defined below.

\begin{align}
\text{Delivery Success Rate} &= \frac{\text{Successfully delivered messages}}{\text{Messages attempted}}, \\
\text{Retransmission Ratio} &= \frac{\text{Retransmitted packets}}{\text{Original packets}}, \\
\text{Goodput} &= \frac{\text{Useful payload bits delivered}}{T_{\text{total}}}.
\end{align}

These metrics give a balanced view across correctness and efficiency.

\section{A Practical Theorem-like Observation}
\begin{theorem}
For a fixed message length and fixed packet loss conditions, increasing sender
window size improves throughput up to a saturation point, after which timeout,
loss recovery, and receiver buffering limits dominate.
\end{theorem}

\begin{Proof}
With small window size, sender frequently waits for acknowledgements and cannot
fully utilize available channel time. Increasing the window allows more
in-flight packets and reduces idle duration. However, once in-flight capacity
matches path-delay product and receiver processing limits, further increase does
not improve delivered payload rate. In lossy conditions, larger windows may also
increase retransmission pressure if recovery logic is not selective. Therefore,
throughput gains plateau after a practical operating point.
\end{Proof}

\clearpage
\section{Interpretation for This Project}
The formulas in this chapter explain observed behavior from level progression:

\begin{enumerate}
    \item versions v0--v2 have very high packet overhead,
    \item v5--v7 improve correctness but remain latency-bound,
    \item v9 and above improve utilization through pipelining,
    \item v10 and above reduce packet count through chunking,
    \item v12 and above improve timing decisions through adaptive RTO,
    \item v13 and above balance aggressiveness and stability using congestion
    control logic and recovery-aware retransmission behavior,
    \item v15--v24 extend the same transport ideas with richer TCP-like
    refinements, including timestamp and recovery-oriented control.
\end{enumerate}

This quantitative interpretation supports the implementation narrative in the
previous chapter and sets context for the comparative results in the next
chapter.